\label{Conclusions}

This project asked whether reinforcement learning can design adaptive
quantitative MRI sequences that estimate tissue relaxation more efficiently
than fixed clinical protocols, and built the full simulation-in-the-loop
pipeline needed to test that question honestly. This chapter draws the three
technical chapters together: what was achieved against the challenges set out
in Chapter~\ref{ch1}, the limitations that bound those achievements, and the
work that would take the project further.

\section{Achievements}
\label{sec:concl-achievements}

The three challenges (C1, simulator validation for qMRI; C2, computational
efficiency under expensive simulation; C3, adaptive sequence design from
tissue-parameter estimates) map onto three achievements, one per technical
chapter.

\textbf{A1 --- a digital twin of the calibration phantom
(Chapter~\ref{ch:mri-system-phantom}).} \texttt{MRISystemPhantom.jl} is a
configurable, open-source digital twin of the QalibreMD Model 130 system
phantom that returns standard \texttt{KomaMRI.Phantom} objects, ships a
tested forward pipeline (phantom construction, sequence generation,
reconstruction, ROI extraction, and fitting), and exposes the randomisation
hooks the RL loop needs. It is the ground-truth foundation for C1 and a
reusable contribution in its own right: a digital twin of a widely deployed
calibration device against which simulated qMRI methods can be checked
before scanner deployment.

\textbf{A2 --- a validated Bloch-simulation pipeline
(Chapter~\ref{ch:validating-komamri}).} Validating that pipeline against the
phantom's known $T_1$ values --- the strict test that distinguishes qMRI from
plausible imaging --- exposed two floating-point time-discretisation failures
in KomaMRI that only appear at the long cumulative sequence times RL training
demands. Both were reduced to minimal reproducers, traced to specific
$\epsilon$-scale timing mechanisms, and fixed upstream; the no-noise
14-sphere $T_1$ fit improved from 39.4\% to 0.48\% mean MAPE. Together with
A1 this answers C1, and the fixes are a contribution back to the wider
simulation community, not only to this project.

\textbf{A3 --- adaptive qMRI sequence design via multi-fidelity RL
(Chapter~\ref{ch:adaptive-rl}).} A full-Bloch step at the target
configuration measured ${\approx}4$ s on CPU, making a single-fidelity run
roughly nine days before evaluation. The multi-fidelity curriculum answers
C2 by learning cheaply on an analytic surrogate and cached or coarsened
water, then warm-starting up a fidelity ladder spanning a ${\sim}30\times$
per-step cost range, with promotion gated on held-out \emph{full-simulator}
validation so a biased rung cannot certify its own policy. On the adaptive
question (C3), the learned closed-loop policy beat the matched fixed log-grid
schedule at two scan budgets on strict held-out episodes (4.62\% versus
6.86\% MAPE at 240 s; 4.16\% versus 6.04\% at 560 s), with non-overlapping
confidence intervals. A memory ablation then produced the project's strongest
result: giving the policy an order-invariant histogram of the inversion times
it has already executed reached 2.93\% MAPE, with 95.8\% of episodes under
5\% --- better than the memoryless control, a fitter-uncertainty channel, a
recurrent policy, and every fixed schedule, on every active sphere. The value of the next acquisition depends
on what has already been measured, and exposing that history explicitly
matters. Translated into the acceleration metric of the closest published
adaptive scheme (Beracha et al.~\cite{adaptive_mri:23}), this
$6.04\%\!\to\!2.93\%$ precision gain corresponds to a nominal ${\approx}4\times$
acceleration --- the top of their reported $1.7$--$4\times$ simulated-$T_1$
range, reached while generalising from their single voxel to a five-sphere
fleet; for $T_1$ both schemes stay in simulation, their hardware validation
covering only $T_2$/MRS.

\section{Limitations}
\label{sec:concl-limitations}

\textbf{Digital twin (Chapter~\ref{ch:mri-system-phantom}).} The twin models
the three contrast plates and fiducial grid but not the resolution insets or
slice-profile wedges, and the sphere-ring geometry matches the manual
photographs rather than a measured device. Two material entries are
approximate: PD-plate relaxation values use bulk water, and the fiducial
relaxation values are generic placeholders. Reference values are fixed at
$20\,^{\circ}$C, with the manual's temperature corrections unimplemented.
None of these affect the $T_1$- and $T_2$-plate experiments reported here,
but all must be addressed before any simulation-to-real comparison involving
the PD plate or fiducials. Coarsening the background water also perturbs the
boundary geometry, a bias quantified rather than eliminated.

\textbf{Simulator validation (Chapter~\ref{ch:validating-komamri}).} The
validation covers the IR-SE sequence family and an idealised scanner: RF
pulses are exact and field inhomogeneities and eddy currents are not modelled.
Other sequence families would need the same timing and no-noise recovery
checks before being trusted for quantitative claims, and the second timing
fix (PR~\#789) remains open upstream, pinned in a project fork.

\textbf{Adaptive RL (Chapter~\ref{ch:adaptive-rl}).} The positive result is a
proof of concept on a deliberately simplified task, not solved full-phantom
mapping: on the full 14-sphere plate the agent learned a genuine
state-dependent policy but still lost to robust fixed coverage, because that
task rewards global experimental design over per-episode adaptation. The
evaluations rest on 24 episodes and one training seed per arm, so they do not
bound sensitivity to initialisation or curriculum timing; the recurrent-policy
negative is further confounded by its ${\approx}5\times$ per-step cost under a
shared wall-clock budget. The action space is narrow (one IR-SE block, mostly
TI and TR), the agent acts only on fitted estimates rather than images, and
every result is simulator-only at a single 1\,mm slab --- the planned MR-Linac
hardware validation was not reached.

\section{Future work}
\label{sec:concl-future}

The most immediate extensions are concrete. For the twin
(Chapter~\ref{ch:mri-system-phantom}), a boundary-aware water grid, verified
PD/fiducial values, and temperature correction of the relaxation tables would
remove the known biases and unlock the PD plate and any sim-to-real
comparison; modelling the omitted resolution insets and slice-profile wedges
would extend it to geometric-distortion and slice-profile experiments. The
twin's architecture also already separates labelled material descriptors from
voxelisation, so other calibration phantoms --- including the open-hardware
designs of Section~\ref{sec:phantom-related-tools} --- could be supported by
writing new descriptor builders while reusing the voxelisation, randomisation,
reconstruction, and fitting pipeline, turning the single-phantom twin into a
small framework for executable qMRI phantom twins. For the simulator
validation (Chapter~\ref{ch:validating-komamri}), the parameter-recovery test
should be repeated for the other shipped sequence families (spin echo,
multi-echo, turbo spin echo, spoiled GRE) before any of them carries
quantitative claims, the still-open PR~\#789 should be landed upstream, and
the idealised-scanner assumption should be relaxed to include RF
inhomogeneity, eddy currents, and $B_0$ imperfections. For the RL environment
(Chapter~\ref{ch:adaptive-rl}), a step-matched recurrent-policy rerun would
separate the memory mechanism from its compute cost, and the fidelity ladder
could be leaned out --- water coarsening proved less useful than caching, so
dropping that rung and re-testing against the biased \texttt{full3} stage is
the obvious simplification.

The scientifically important direction is to widen the action space and move
beyond the $T_1$ plate. Letting the agent choose the acquisition family
(IR-SE, spin echo, multi-echo trains, spoiled GRE) and the relevant per-family
parameters would test the more publishable claim that RL's advantage is in
choosing \emph{which contrast mechanism} to deploy, and would extend naturally
to joint $T_1/T_2/T_2^*$ and off-resonance estimation --- closer to magnetic
resonance fingerprinting. Two realism steps accompany this: replacing the
hard-pulse approximation with shaped slice-selective (windowed-sinc) pulses,
and generalising the single thin slab to genuine 3-D slice selection, both of
which the twin already supports but the experiments do not yet exercise.
Replacing the single fixed scan budget with an accuracy--time Pareto objective,
via a stop action or scan-time penalty, would reflect the real clinical
trade-off, where reaching adequate accuracy sooner can matter more than a lower
asymptotic error. Finally, executing a learned protocol on the MR-Linac would
close the sim-to-real gap the project was ultimately motivated by, testing
whether the learned timings survive real scanner constraints, reconstruction,
and phantom data.

\section{Concluding remarks}
\label{sec:concl-closing}

The project delivered a validated, reusable digital twin; surfaced and fixed
real numerical bugs in a widely used simulator; and built a bias-aware
multi-fidelity training loop that made simulation-in-the-loop qMRI RL
tractable. Within that pipeline it showed, on a controlled task, that an
adaptive policy with an explicit memory of its own acquisitions can beat
strong fixed protocols. To our knowledge this is the first reinforcement-learning
agent for adaptive quantitative MRI: a policy that conditions each acquisition
on the current fitted tissue parameters and is trained and evaluated directly
on fitted-$T_1$ error, where prior adaptive qMRI derives its acquisition rule
from a Bayesian model and prior RL in MRI targets non-quantitative objectives
such as shape classification or k-space sampling. The honest boundary of the
work is equally clear: adaptivity paid only where the task was constructed to
require it, and every result remains simulated.

The wider context is the one set out in Chapter~\ref{ch1}: MRI demand is
growing faster than capacity, and adaptive acquisition attacks the scan-time
bottleneck at the measurement-design level rather than after the data are
collected. Within MRI digital-twin research, which remains concentrated on
diagnosis and treatment planning, protocol- and infrastructure-level twins
are the underdeveloped fraction, held back by simulation cost and
validation~\cite{greggio_digital_twinning_mri_2026} --- precisely the two
barriers the multi-fidelity ladder and the parameter-recovery validation
address in this bounded calibration setting. The phantom twin, validated
simulator, and training loop together form a testbed in which adaptive qMRI
protocols can be developed and compared before any patient or scanner is
involved. Turning this controlled demonstration into a robust,
hardware-validated qMRI sequence designer is the natural continuation, and
the infrastructure built here is intended to make that continuation
straightforward.
